% =====================================================================================
% BORRADOR — Sección del prototipo (sec:prototipo) y su evaluación
% Generado el 2026-10-09 a partir de docs/ del repositorio control-acceso-simple.
% Fuente de autoridad: docs/guia-tesis.md (qué afirmar y qué no).
%
% Paquetes que necesita en el preámbulo de la tesis (si no están ya):
%   \usepackage{booktabs}            % tablas
%   \usepackage{tikz}                % figura de arquitectura
%   \usetikzlibrary{arrows.meta,positioning}
%   \usepackage{url} o hyperref      % \url
%
% Citas: solo se usan claves que existen en docs/referencias-comparacion-rostros.bib.
% Lo que falta se marca \cite{PENDIENTE-<tema>} y está listado al final de este archivo.
%
% Ubicación sugerida (docs/guia-tesis.md):
%   - \section{...}\label{sec:prototipo} ............ capítulo de Metodología (guía §2 y §3)
%   - \section{...}\label{sec:prototipo-evaluacion} . capítulo de Resultados (guía §4, §5, §6)
%   - \section{...}\label{sec:prototipo-limitaciones} Discusión o Conclusiones (guía §8)
% =====================================================================================

% --- Cifras que se están volviendo a medir ------------------------------------------
% TODO: el experimento biométrico se está repitiendo con el umbral FAR corregido (la
% medición anterior aceptaba un impostor de más). Reemplazar estos valores con los de
% docs/results/biometric-degradation.md (fila 112x112@900B y línea base) cuando termine.
\newcommand{\TARqr}{\textbf{[pendiente]}}         % TAR @ FAR 1e-3, foto del QR 112x112 en 900 B
\newcommand{\TARplantilla}{\textbf{[pendiente]}}  % TAR @ FAR 1e-3, plantilla del emisor (línea base)
\newcommand{\UmbralBio}{\textbf{[pendiente]}}     % umbral coseno que usa el escáner (biometrics.ts)

% --- Cifras de verificación del software ---------------------------------------------
% TODO: actualizar al cerrar la tesis con `pytest` (backend) y `npm run test` (frontend).
\newcommand{\TestsBackend}{138}
\newcommand{\TestsFrontend}{85}

% =====================================================================================
\section{Prototipo de verificación por código QR dinámico}
\label{sec:prototipo}
% =====================================================================================

Esta sección describe el prototipo de verificación por código QR dinámico construido para esta
tesis. Es un componente independiente del bosquejo conceptual del objetivo específico~5: lo
demuestra de forma aplicada, pero no lo reemplaza. El caso de aplicación es el control de acceso
a eventos masivos (estadios, conciertos, festivales), donde hoy se lee el código QR estático del
carnet físico y el personal compara a simple vista la fotografía impresa con la persona.

El prototipo \emph{simula} la aplicación Identificación Digital del Servicio de Registro Civil e
Identificación: el enrolamiento, el PIN y el QR dinámico. Lo hace porque las claves y las
interfaces de verificación del sistema real no son públicas; por lo mismo, no interopera con él ni
describe su funcionamiento interno. El aporte que se evalúa es de \textbf{seguridad}, no de
velocidad: la hipótesis es que el QR dinámico aumenta la seguridad de la verificación sin alargar
de forma relevante las filas. El código, las decisiones de diseño con su justificación y los
experimentos reproducibles están en el repositorio del prototipo\footnote{%
  \url{https://github.com/mmonzodiaz/TESIS---Control-QR-Dinamico}. % TODO: confirmar que el repositorio es público o cambiar por anexo.
}; aquí se resumen.

% -------------------------------------------------------------------------------------
\subsection{Actores y arquitectura}
\label{sec:prototipo-arquitectura}

El diseño sigue el modelo de tres actores de las credenciales verificables
\cite{PENDIENTE-w3c-vc}: un \emph{emisor} que firma la credencial, un \emph{titular} que la porta y
un \emph{verificador} que la comprueba. A ellos se suma el sistema del recinto, que conoce las
entradas vendidas pero nunca recibe la fotografía de los asistentes
(figura~\ref{fig:prototipo-arquitectura}).

\begin{itemize}
  \item \textbf{Registro Civil simulado (emisor).} Mantiene un padrón de personas, emite la cédula
        digital firmada y publica la lista firmada de cédulas revocadas. Es un servicio separado del
        recinto, que solo conoce su clave pública.
  \item \textbf{Teléfono del titular.} Guarda la cédula cifrada con un PIN y genera el QR dinámico
        sin conexión.
  \item \textbf{Escáner del guardia (verificador).} Descifra y verifica el QR localmente, muestra al
        guardia la fotografía firmada y consulta al recinto sin enviarle la fotografía.
  \item \textbf{Sistema del recinto.} Tiene la lista de entradas nominativas por RUT, las
        prohibiciones de ingreso y el historial; decide uso único, entrada y prohibiciones, y
        entrega al escáner la clave de lectores autorizados y una copia para operar sin red.
\end{itemize}

\begin{figure}[htbp]
  \centering
  \begin{tikzpicture}[
      actor/.style={draw, rounded corners, align=center, minimum width=4.2cm, minimum height=1.3cm},
      flujo/.style={-{Stealth[length=2.5mm]}, thick},
      etiqueta/.style={font=\footnotesize, align=center}
    ]
    \node[actor] (emisor)   {Registro Civil simulado\\\footnotesize(emisor)};
    \node[actor, right=5.2cm of emisor] (recinto) {Sistema del recinto\\\footnotesize(entradas, prohibiciones, historial)};
    \node[actor, below=3.2cm of emisor] (telefono) {Teléfono del titular\\\footnotesize(cédula cifrada con PIN)};
    \node[actor, below=3.2cm of recinto] (escaner) {Escáner del guardia\\\footnotesize(verificador)};

    \draw[flujo] (emisor) -- node[etiqueta, left] {cédula firmada\\(Ed25519)} (telefono);
    \draw[flujo] (telefono) -- node[etiqueta, below] {QR dinámico cifrado\\(generado sin red)} (escaner);
    \draw[flujo] (emisor) -- node[etiqueta, above] {lista de revocadas firmada} (recinto);
    \draw[flujo] ([xshift=-6mm]escaner.north) -- node[etiqueta, left] {consulta\\sin la foto} ([xshift=-6mm]recinto.south);
    \draw[flujo] ([xshift=6mm]recinto.south) -- node[etiqueta, right] {veredicto; copia\\sellada para\\operar sin red} ([xshift=6mm]escaner.north);
  \end{tikzpicture}
  \caption{Actores del prototipo. La fotografía viaja firmada dentro del QR, del teléfono al escáner,
    y nunca llega al sistema del recinto. El recinto entrega al escáner autenticado la clave de
    lectores autorizados y una copia de contingencia sellada.}
  \label{fig:prototipo-arquitectura}
\end{figure}

% -------------------------------------------------------------------------------------
\subsection{La cédula digital}
\label{sec:prototipo-cedula}

La cédula es de la persona y no del evento: sirve ante cualquier lector autorizado, y el recinto
solo coteja el RUT firmado contra su lista de entradas. Contiene el RUT, el nombre, el resumen
SHA-256 de la fotografía y la clave pública Ed25519 del teléfono en que se activó, todo firmado
por el emisor con Ed25519 \cite{PENDIENTE-rfc8032}. Que la clave del teléfono vaya dentro de lo
firmado liga la cédula a ese dispositivo: una copia del archivo en otro teléfono no sirve, porque
no puede producir la firma del dispositivo que exige cada QR.

Para activarla, el titular ingresa su RUT y una clave de activación de ocho caracteres que el
emisor tiene en su padrón; luego define un PIN de seis dígitos que cifra la cédula y la clave del
dispositivo en el teléfono (cinco fallos seguidos bloquean el ingreso por 60~s). Esta clave de
activación es un \textbf{sustituto simulado}: la aplicación real activa con ClaveÚnica más
verificación facial con prueba de vida \cite{PENDIENTE-activacion-app}. Cada persona tiene una sola
cédula activa: activarla en otro teléfono revoca la anterior.

% -------------------------------------------------------------------------------------
\subsection{El código QR dinámico}
\label{sec:prototipo-qr}

El QR no lo emite un servidor: lo genera el teléfono, sin conexión, cada 45~s. Contiene la cédula
firmada, la fotografía y una firma del dispositivo sobre el resumen de la cédula, un número
aleatorio de un solo uso (\emph{nonce}) y la hora de emisión, con separación de dominio entre la
firma del emisor y la del dispositivo. Todo va cifrado para la clave pública de los lectores
autorizados mediante un acuerdo X25519 con clave efímera \cite{PENDIENTE-rfc7748}, derivación HKDF
\cite{PENDIENTE-rfc5869} y cifrado autenticado ChaCha20-Poly1305 \cite{PENDIENTE-rfc8439}, y se
codifica en base45 \cite{PENDIENTE-rfc9285} para aprovechar el modo alfanumérico del QR. La
tabla~\ref{tab:prototipo-qr} resume qué garantiza cada pieza.

\begin{table}[htbp]
  \centering
  \small
  \begin{tabular}{p{4.2cm}p{9.3cm}}
    \toprule
    Mecanismo & Qué impide \\
    \midrule
    Firma del emisor (Ed25519) & Fabricar una cédula, o cambiar la fotografía, el RUT o el nombre de una real. \\
    Firma del dispositivo (Ed25519) & Generar QR con una cédula copiada a otro teléfono. \\
    Hora firmada y vigencia de 45~s & Reutilizar una fotografía o captura del QR más tarde. \\
    \emph{Nonce} de un solo uso & Que un mismo QR sirva para más de un ingreso. \\
    Cifrado para lectores autorizados & Que quien fotografía el QR en la fila lea la identidad o la fotografía. \\
    \bottomrule
  \end{tabular}
  \caption{Mecanismos del QR dinámico del prototipo y lo que cada uno impide.}
  \label{tab:prototipo-qr}
\end{table}

La fotografía se comprime a $112\times112$~píxeles en un máximo de 900~bytes, lo que produce un
QR denso (unos 1{,}7~kB) que se leyó sin dificultad con la cámara de un teléfono y con una pistola
lectora 2D en una primera prueba con hardware real. Como contraste, la caracterización externa del
QR de la aplicación oficial (escaneo del código de la cédula digital del autor con un lector
genérico, sin decodificar su contenido) mide 256~bytes exactos, donde no cabe una fotografía; su
verificador necesita otro canal. Entregar al guardia una fotografía autenticada verificable
\emph{sin conexión} es, por tanto, una decisión de diseño de este prototipo y no una
característica heredada de la aplicación real. % Guía §1.1(b): no llamarlo «ingeniería inversa».

\subsubsection{Vigencia de 45 segundos}
\label{sec:prototipo-ttl}

La vigencia no es arbitraria: se ubica entre dos cotas. La inferior es de usabilidad: un escaneo
legítimo toma entre 5 y 10~s (mostrar la pantalla, enfocar, consultar al servidor), y con un
margen de 2 a 4 veces para los arquetipos con menor familiaridad digital de la encuesta, el piso
queda entre 15 y 30~s. La superior es de exposición a reutilización: el precedente formal más
cercano es el paso de 30~s de los códigos de un solo uso basados en tiempo \cite{PENDIENTE-rfc6238},
y los esquemas de pago por QR dinámico en uso rotan en ventanas de decenas de segundos. Se fijó
45~s como punto intermedio, con una tolerancia de reloj de 5~s: una copia expone a lo más 50~s.
No existe un valor público del intervalo de la aplicación oficial; el prototipo replica su
propósito declarado (impedir la reutilización por captura), no un parámetro suyo.

Como el teléfono genera los QR sin conexión, dentro de esa ventana pueden coexistir varios QR
vigentes de la misma cédula. Lo que acota el efecto es la admisión única por entrada
(sección~\ref{sec:prototipo-verificacion}) y que una copia usada por un tercero muestra al guardia
la fotografía del titular.

% -------------------------------------------------------------------------------------
\subsection{Verificación en la puerta}
\label{sec:prototipo-verificacion}

El escáner descifra el QR, verifica ambas firmas y la vigencia, y muestra de inmediato la
fotografía firmada en la pantalla del guardia. En paralelo consulta al sistema del recinto
\emph{sin enviarle la fotografía}, en dos pasos:

\begin{enumerate}
  \item \textbf{Verificación} (\texttt{check}): el recinto revisa firmas, revocación, entrada
        nominativa del RUT para el evento de esa fila, uso único, ingreso previo y prohibiciones, sin
        consumir nada ni registrar el ingreso.
  \item \textbf{Decisión del guardia}: si la fotografía corresponde a quien tiene al frente, confirma
        (\texttt{admit}) y se registra el ingreso; si no, rechaza (\texttt{reject}). El rechazo
        inutiliza ese QR pero \emph{no} consume la entrada, de modo que el titular legítimo puede
        entrar después con un QR nuevo. La confirmación se acepta hasta 60~s después de la
        verificación, aunque el QR haya vencido mientras el guardia miraba la foto.
\end{enumerate}

El uso único es atómico aunque varias filas escaneen a la vez: el \emph{nonce} se inserta como
clave primaria y solo una inserción puede ganar. Una persona ya admitida queda bloqueada con
cualquier QR nuevo en ese evento (anti \emph{pass-back}), y las prohibiciones de ingreso se revisan
en cada validación para que una sanción agregada a última hora aplique.

La diferencia con mostrar la fotografía en la pantalla del titular es central: una aplicación
falsa o un video pueden mostrar cualquier imagen; lo que tiene valor es la fotografía que el lector
del guardia obtiene del QR firmado.

% -------------------------------------------------------------------------------------
\subsection{Revocación}
\label{sec:prototipo-revocacion}

Activar la cédula en otro teléfono o reportar la pérdida («¿Perdiste tu teléfono?», con el RUT y
la clave de activación) revoca la emisión anterior. El emisor publica una lista firmada con los
resúmenes de las emisiones revocadas, sin RUT ni nombres. El recinto la consulta como máximo cada
30~s y rechaza con \texttt{revoked}; si el emisor no responde, sigue con la última lista válida y,
si nunca la obtuvo, el escáner se lo advierte al guardia en vez de bloquear a todos. No es un
servicio de estado estandarizado (\emph{status list}, OCSP); es una simplificación declarada.

% -------------------------------------------------------------------------------------
\subsection{Operación sin conexión}
\label{sec:prototipo-contingencia}

En un estadio la red puede caerse. Mientras hay conexión, el escáner baja cada minuto una copia de
contingencia del evento (RUT con entrada, ya admitidos, prohibiciones y revocadas) que guarda solo
en memoria. Sin red, verifica igual el QR y la fotografía, rechaza según la copia y deja decidir
al guardia; con una copia de más de 10 minutos no deja pasar a nadie y deriva a la fila normal.

Las decisiones tomadas sin red quedan en una cola del lector, separada por guardia, y se
sincronizan al volver la conexión. El recinto las reverifica como si se hubieran validado a la hora
del escaneo y marca como \emph{alerta} todo ingreso que no habría permitido. Para que esa hora no
se pueda falsear, el servidor sella cada copia con un HMAC sobre el evento y su hora de
generación, y solo acepta la hora declarada dentro de la vigencia de la copia sellada con que se
decidió, o de la ventana de un \texttt{check} en línea previo; sin ninguno de los dos, evalúa a la
hora de la sincronización. % TODO: si se incorpora el tope de antigüedad del sello (revisión del 2026-10-09, N3), mencionarlo aquí.
Una alerta no consume la entrada del titular: quien entró con un QR vencido o ya usado
probablemente no era el titular.

El riesgo residual se declara: dos puertas sin red no se enteran entre sí, así que la misma persona
puede ingresar por ambas; el sistema lo \emph{detecta} al sincronizar, pero no lo \emph{evita}.

% -------------------------------------------------------------------------------------
\subsection{Comparación biométrica como módulo opcional}
\label{sec:prototipo-biometria}

A pedido de la revisión, el prototipo incorpora un comparador facial que contrasta la fotografía
firmada con la imagen de la cámara del guardia. Está \textbf{implementado y medido en laboratorio,
pero no validado en terreno}, y viene \textbf{apagado por defecto}. Sus invariantes son:
\begin{itemize}
  \item corre en el dispositivo del guardia, nunca en el teléfono del titular (que podría informar
        un resultado falso) ni en el servidor;
  \item no guarda imágenes ni vectores faciales: el puntaje se muestra al guardia y se descarta;
  \item el puntaje informa, pero decide el guardia con los mismos botones de la verificación en dos
        pasos;
  \item siempre existe una vía de ingreso no biométrica.
\end{itemize}

El último invariante tiene fundamento legal. El rostro es dato personal sensible y, para tratarlo,
la única base de licitud aplicable es el consentimiento expreso; el art.~12 inc.~6° de la
Ley~N°~19.628, en el texto fijado por el art.~primero de la Ley~N°~21.719, presume que el
consentimiento no fue libre si la recolección no era necesaria para el servicio
\cite{PENDIENTE-ley21719}. Si la comparación fuera el único modo de entrar, esa presunción haría
caer la base de licitud. La ley entra en vigencia el 1 de diciembre de 2026, con un proyecto de
prórroga en trámite: el prototipo se diseña contra la norma que regiría al desplegarlo.
% TODO: reverificar el estado de la prórroga (Boletín 18.623-07) al cerrar la tesis.

El costo de llevar la fotografía dentro del QR se midió con 6000 pares de rostros del conjunto LFW
\cite{PENDIENTE-lfw} y un modelo de la familia ArcFace \cite{PENDIENTE-arcface}: comparar contra la
foto de $112\times112$ en 900~bytes alcanza una tasa de aceptación genuina (TAR) de \TARqr{} a una
tasa de falsa aceptación (FAR) de $10^{-3}$, frente a \TARplantilla{} de una plantilla calculada por
el emisor desde el retrato original; el umbral del escáner es \UmbralBio{}. Son cotas de
laboratorio: sin alineado facial, con fotografías de archivo y no con fotogramas de una fila, y
con un conjunto que no representa a la población chilena. Los retratos de la demostración son
sintéticos, y medir el rendimiento real exigiría rostros con consentimiento informado.

Como contraste se construyó también un módulo secundario, apagado por defecto, en que el emisor
entrega la fotografía a partir del núcleo firmado en vez de llevarla en el QR. Gana algunos puntos
de TAR, pero pierde la operación sin conexión, vuelve a poner la fotografía en la red en cada
escaneo y permitiría al recinto reunir la fotografía de cada asistente; por eso no es el flujo
principal.

% -------------------------------------------------------------------------------------
\subsection{Seguridad de la aplicación y verificación del software}
\label{sec:prototipo-software}

Además del protocolo, la aplicación web se endureció siguiendo OWASP Top~10 \cite{PENDIENTE-owasp}:
contraseñas del personal con Argon2id y política de NIST SP~800-63B \cite{PENDIENTE-nist80063b},
sesiones JWT con secreto obligatorio que se invalidan al cambiar la contraseña o el rol, consultas
parametrizadas, cabeceras de seguridad y política de contenido, límites de intentos y registro de
eventos de seguridad. El historial no expone datos personales del personal y enmascara el RUT de los
asistentes.

La implementación está cubierta por \TestsBackend{} pruebas automáticas del servidor y
\TestsFrontend{} del cliente, entre ellas vectores de prueba compartidos que obligan a que ambos
produzcan el mismo QR byte a byte. El 9 de octubre de 2026 se hizo una revisión completa del código
que encontró 34 problemas, entre ellos que la sincronización sin red aceptaba cualquier hora de
escaneo pasada y que el umbral biométrico aceptaba un impostor de más. Todos se corrigieron con una
prueba que falla con el código anterior, y una segunda revisión verificó el cierre.
% TODO (autor): declarar la herramienta con que se hicieron ambas revisiones según la política de
% uso de IA de la universidad, y actualizar el detalle si cambia.

% =====================================================================================
\section{Evaluación de seguridad del prototipo}
\label{sec:prototipo-evaluacion}
% Según docs/guia-tesis.md §4-§6, esta sección va en el capítulo de Resultados.
% =====================================================================================

\subsection{Comparación con el QR estático del carnet}
\label{sec:prototipo-carnet}

El control actual en eventos lee el QR del carnet físico, que contiene una dirección de consulta de
vigencia del Registro Civil con el RUN, el tipo y el número de documento y la línea media de la
zona de lectura mecánica; ninguna fuente oficial especifica ese contenido, que se establece por
triangulación de tres evidencias independientes \cite{PENDIENTE-docstatus}. De ahí se siguen sus
debilidades: es fijo durante toda la vida del documento (una fotografía o fotocopia sirve igual),
no va firmado, no trae la fotografía (la que se compara es la impresa en la tarjeta, que en un
carnet falso controla el falsificador) y expone la fecha de nacimiento a quien lo fotografíe. La
fotografía del carnet va grabada con láser en blanco y negro \cite{srcei_resolucion466_2024}.

Conviene no exagerar el contraste: el carnet \emph{sí} tiene criptografía, pero en el chip y no en
el QR \cite{srcei_resolucion466_2024}; aprovecharla en una puerta exigiría lectores NFC y el
directorio de claves públicas de ICAO. La tabla~\ref{tab:prototipo-carnet} resume la comparación.

\begin{table}[htbp]
  \centering
  \small
  \begin{tabular}{p{3.6cm}p{4.9cm}p{5.0cm}}
    \toprule
    Propiedad & QR del carnet físico & QR dinámico del prototipo \\
    \midrule
    Copia diferida & Sirve indefinidamente & Vence a los 45~s (más 5~s de tolerancia) \\
    Autenticidad & Sin firma: se puede generar o alterar & Firmas del emisor y del dispositivo \\
    Fotografía de referencia & Impresa en la tarjeta, que controla el falsificador & Firmada por el emisor, mostrada en el dispositivo del guardia \\
    Datos expuestos al fotografiarlo & RUN, número de documento y fechas & Ninguno (cifrado para lectores autorizados) \\
    Operación sin conexión & Comparación visual de la tarjeta & Verificación local de firmas y foto; copia de hasta 10~min \\
    \bottomrule
  \end{tabular}
  \caption{QR estático del carnet físico frente al QR dinámico del prototipo.}
  \label{tab:prototipo-carnet}
\end{table}

\subsection{Experimento de ataques}
\label{sec:prototipo-ataques}

Los mismos ataques se ejecutaron contra ambos modos sobre el servidor real, con un reloj
controlado para simular esperas; ambos modos tienen la misma regla de admisión única, de modo que
la única variable es que la credencial sea estática o dinámica. Se distinguen dos niveles: si el
ataque pasa la validación automática, y si pasa además una comparación correcta de la fotografía.
La tabla~\ref{tab:prototipo-ataques} resume los resultados.

\begin{table}[htbp]
  \centering
  \small
  \begin{tabular}{p{6.4cm}cc}
    \toprule
    Nivel & QR del carnet & QR dinámico \\
    \midrule
    Ataques que pasan la validación automática & 4 de 5 & 1 de 5 \\
    Ataques que pasan además la comparación de la foto & 4 de 5 & 0 de 5 \\
    Evasiones del filtro de prohibiciones (automático) & 2 de 2 & 1 de 2 \\
    Evasiones del filtro de prohibiciones (con la foto) & 1 de 2 & 0 de 2 \\
    \bottomrule
  \end{tabular}
  \caption{Experimento de ataques. Los ataques son el uso de una fotografía del QR 60~s y una hora
    después, una copia usada 30 días después en otro evento, el reenvío en tiempo real a un cómplice
    y la copia usada después de que el titular ingresó. En ambos modos el titular legítimo es
    aceptado y la persona sancionada es bloqueada con su propia credencial.}
  \label{tab:prototipo-ataques}
\end{table}

Con el QR dinámico, toda copia diferida se bloquea por vigencia y el titular legítimo puede entrar
después; con el carnet, el atacante entra y el titular queda afuera. La fila de la copia usada en
otro evento se bloquea por la vigencia y no por tratarse de otro evento: el QR no está ligado al
evento, porque la cédula es de la persona, así que un reenvío en tiempo real sirve en cualquier
evento donde el titular tenga entrada. Ese reenvío es el único ataque que pasa la validación
automática del QR dinámico, y lo detiene la fotografía: el cómplice queda frente a un guardia que ve
la fotografía firmada del titular, mientras que con un carnet falso la fotografía es la del propio
atacante.

El nivel «con comparación correcta de la fotografía» es un \emph{límite superior}, no el
resultado esperado. La comparación humana de rostros no familiares falla con frecuencia: comparar
a una persona en vivo con una fotografía ronda el 85~\% de acierto \cite{megreya2008matching};
oficiales de pasaportes aceptaron un 14~\% de fotografías fraudulentas, sin relación con su
experiencia \cite{white2014passport}; los impostores poco frecuentes se pasan por alto más
\cite{papesh2014infrequent}, y las fotografías antiguas o de baja resolución empeoran el acierto
\cite{megreya2013monthsapart,bindemann2013pixelation}. Aun así, la conclusión no cambia de signo:
el QR dinámico elimina las copias y obliga a un impostor \emph{parecido} al titular y coordinado en
tiempo real. Que la fotografía de referencia sea a color podría ayudar al guardia, porque comparar
una imagen en escala de grises con una persona vista a color sesga hacia aceptar al impostor
\cite{bobak2019greyarea} y el color aporta más en imágenes pequeñas \cite{yip2002color}; es una
inferencia desde pares de fotografías, no un resultado medido con personas en vivo.
% Guía §1.1(c): argumento de apoyo, no titular. NO afirmar que la app oficial muestra la foto a
% color: la Res. Ex. 234/2026 fija la «fotografía principal en blanco y negro» (docs/fuentes-pendientes.md).

\subsection{Costo en la fila}
\label{sec:prototipo-fila}

La velocidad no es el argumento, pero hay que mostrar que la validación en línea no la empeora. En
la prueba de carga, la consulta al servidor suma unos 24~ms por persona y no hubo ninguna doble
concesión de ingreso; en la corrida publicada, con 16 escáneres simultáneos, uno de 400 titulares
válidos no fue admitido por la espera de la base de datos, un problema de disponibilidad y no de
seguridad. Un modelo de colas M/G/1 validado por simulación estima un aumento de la espera de un
3~\% en la punta. El tiempo humano de servicio (8~s) y la tasa de llegadas son supuestos no
medidos.
% TODO: revisar estas cifras contra docs/queue-model.md antes de integrarlas.

% =====================================================================================
\section{Limitaciones del prototipo}
\label{sec:prototipo-limitaciones}
% Según docs/guia-tesis.md §8, va en Discusión o Conclusiones.
% =====================================================================================

Las limitaciones se declaran para no inflar las garantías:
\begin{itemize}
  \item \textbf{Emisor simulado:} una sola clave, sin certificados ni cadena de confianza.
  \item \textbf{Enrolamiento:} la clave de activación sustituye a ClaveÚnica más verificación facial
        con prueba de vida.
  \item \textbf{Biometría:} módulo opcional apagado por defecto, medido en laboratorio y no validado
        en terreno; sin detección de ataques de presentación es derrotable con una fotografía del
        titular.
  \item \textbf{Sin conexión:} el doble ingreso por dos puertas sin red se detecta, no se evita.
  \item \textbf{Almacenamiento:} la cédula y la clave del teléfono van cifradas con el PIN en el
        navegador, no en un almacén seguro de hardware.
  \item \textbf{Despliegue:} requiere TLS; los límites de intentos viven en la memoria de un solo
        proceso, el registro de seguridad no está centralizado y el personal no tiene segundo factor.
  \item \textbf{Validez legal:} no es firma electrónica ni aporta no repudio.
  \item \textbf{Comparación humana:} el guardia puede equivocarse al comparar la fotografía; los
        resultados con comparación correcta son un límite superior.
\end{itemize}

% =====================================================================================
% CITAS PENDIENTES (no existen en referencias-comparacion-rostros.bib ni en fuentes-pendientes.md)
%   PENDIENTE-w3c-vc        W3C Verifiable Credentials Data Model (modelo emisor-titular-verificador)
%   PENDIENTE-rfc8032       RFC 8032, Ed25519
%   PENDIENTE-rfc7748       RFC 7748, X25519
%   PENDIENTE-rfc5869       RFC 5869, HKDF
%   PENDIENTE-rfc8439       RFC 8439, ChaCha20-Poly1305
%   PENDIENTE-rfc9285       RFC 9285, base45
%   PENDIENTE-rfc6238       RFC 6238, TOTP (paso de 30 s)
%   PENDIENTE-ley21719      Ley N° 19.628, texto fijado por el art. primero de la Ley N° 21.719 (y BCN)
%   PENDIENTE-lfw           Labeled Faces in the Wild (Huang et al., 2007) + variante y espejo usados
%   PENDIENTE-arcface       ArcFace (Deng et al., 2019) / InsightFace w600k_mbf y licencia de sus pesos
%   PENDIENTE-owasp         OWASP Top 10 (2021)
%   PENDIENTE-nist80063b    NIST SP 800-63B
%   PENDIENTE-activacion-app Activación real con ClaveÚnica + verificación facial con prueba de vida
%                           (hoy no hay fuente en el repo; revisar la Res. Ex. 234/2026, srcei_resolucion234_2026)
%   PENDIENTE-docstatus     Formato del QR del carnet (ver docs/threat-model.md §6 y docs/fuentes-pendientes.md:
%                           kliashchou_regula_chile_2026 es secundaria; la verificación propia va como metodología)
% Antes de entregar: pasar todas las citas por la skill verificador-citas.
% =====================================================================================
